\documentclass[11pt]{article}

\usepackage{amsmath,amssymb,amstext,amsthm}
\usepackage{geometry}
\usepackage{fancyhdr}
\usepackage[pdftex]{graphicx}
\usepackage{xcolor}

\geometry{
  verbose,
  dvips,
  width=430pt, marginparsep=0pt, marginparwidth=0pt,
  top=90pt, headheight=12pt, headsep=20pt, footskip=30pt, bottom=80pt}

\setlength{\parskip}{\medskipamount}
\setlength{\parindent}{0pt}

\linespread{1.05}

\newtheorem{theorem}{Theorem}
\newtheorem{lemma}[theorem]{Lemma}
\newtheorem{proposition}[theorem]{Proposition}
\newtheorem{corollary}[theorem]{Corollary}
\newtheorem{conjecture}[theorem]{Conjecture}
\newtheorem{obs}{Observation}
\newtheorem{clm}{Claim}
\newtheorem{ques}{Question}
\newtheorem{problem}{Problem}

\newtheorem{ex}{Example}


\newcommand{\spc}{\hspace{0.08in}}
\newcommand{\vs}{\vspace*{.1in}}
\newcommand{\E}{\mathbb{E}}
\newcommand{\Prob}{\mathbb{P}}
\newcommand{\edit}[1]{\emph{\textcolor{blue}{#1}}}
\newcommand*\dif{\mathop{}\!\mathrm{d}}


\renewenvironment{proof}{{\noindent \textbf{\textit{Proof.}}}}
{\hfill $\Box$ \vspace*{0.1in}}
\newenvironment{proofc}{{\noindent \textit{Proof of Claim.}}}
{\hfill $\Box$}



\begin{document}

\begin{LARGE}\begin{center}\bf 16.2 Line Integrals\end{center}
\end{LARGE}

\vs\vs



\section{Warm Up}
\textbf{Question 1:} Find the arc length of the vector function $\vec{r}(t) = \hat{i}+\hat{j}$ from $t = 1$ to $t = 5$.

\section{Motivation}
When we see $\int_a^b f(x) \dif{x}$ we often think of area under the curve $f(x)$ bounded by the interval $[a,b]$. While this is true, we also know from our chapter 15 that we can also think of this as the mass of a thin wire from $a$ to $b$ where $f(x)$ is the mass density function of the wire. We then extended that to finding the mass of a region with $\iint_R f(x,y) \dif{A}$ and the mass of a solid with $\iiint_T f(x,y,z) \dif{V}$. Once again, $f(x,y)$ and $f(x,y,z)$ are mass density functions. Our goal for today is to define line integrals when our thin wire is no longer straight. How would I use integrals to find the mass of a wire with mass density function.



\section{Definitions \& Theorems}

\begin{enumerate}
\item If we have some curve $C$ and we want to know the mass of the curve we can use a line integral. If our curve is not just a straight line, it would be in terms of two variables $x$ and $y$ in other words, it would start at some point $(x_1,y_1)$ and end at some terminal point $(x_2,y_2)$. The issue is that mass of a wire is single integral but we have a curve which is of two variables. How do we fix this?

\item The other way to imagine this same problem is if I had a curve $C$ projected on some surface $f(x,y)$. It is important to note that the curve is a 2D curve but it can be projected into 3-space. The line integral will find the area under this curve projection. Numerically this will be equal to the mass of the curve.

\item We also no longer have a $\dif{x}$ anymore since our region is not an interval along the $x$-axis. Instead we have to remember back from vector functions and realize that it is arc length which we now are adding up infinity small pieces of to get my entire curve $C$! we use $\dif{s}$ now for arc length.

\item To define the line integral mathematically we need to define our curve $C$ (which isn't a function) using parametric equations. So now $C = \vec{r}(t)$  We get the following.

\begin{align*}
   \text{Mass of wire} &= \int_C f(x,y) \dif{s}\\
   &= \int_{t=a}^{t=b} f\left(x(t),y(t)\right) \left|\vec{r'}(t)\right| \dif{t} 
\end{align*}

\item Once again the line integral tells us the area under the surface $f(x,y)$ bounded by a curve $C$ (which is of 2 variables) and this number is numerically equivalent to the mass of the curve $C$ when $f(x,y)$ is the mass density function of $C$. 

\item Wait a minute... We learned that $\iint_R 1 \dif{A}$ was equal to the area of the region $R$ and $\iiint_T 1 \dif{V}$ was equal to the volume of $T$. So triple integral of $1$ gives volume, double integral of $1$ gives area and so the single integral of $1$ gives length. We already know this from chapter 13!!!

\begin{equation*}
    \int_{t=a}^{t=b} |\vec{r'}(t)| \dif{t} = \text{arc length!}
\end{equation*}

\item We now are going to imagine our space curve traveling through a vector field. We can utilize the line integral to represent the amount of work being done as a particle travels along the curve through the vector field. To compute this we use the following formula.

\begin{equation*}
    w = \int_C F(\vec{r}(t)) \cdot \vec{r'}(t) \dif{t} = \int_C F \cdot \dif{r}
\end{equation*}

\end{enumerate}




\section{Examples}

\begin{problem}
    Evaluate the following line integral, where $C: \vec{r} = 2t\hat{i}+t^3\hat{j}$, $0\leq t\leq 1$.

    \begin{equation*}
        \int_C y \dif{s}
    \end{equation*}
\end{problem}

\vspace{8cm}


\begin{problem}
    Evaluate the following line integral, where $C$ is the right half of the circle $x^2+y^2=9$.

    \begin{equation*}
        \int_C (x^2+y^2) \dif{s}
    \end{equation*}
\end{problem}

\vspace{7cm}


\begin{problem}
    Evaluate the following line integral where $C$ is the line segment from $(-2,1)$ to $(1,3)$.

    \begin{equation*}
        \int_C 2xy \dif{s}
    \end{equation*}
\end{problem}

\vspace{9cm}

\begin{problem}
    Evaluate the following line integral where $C$ is the line segment from $(1,1,0)$ to $(2,3,1)$.

    \begin{equation*}
        \int_C xyz^2 \dif{s}
    \end{equation*}
\end{problem}


\vspace{9cm}

\begin{problem}
    Find the work done by moving along the parabola $y=x^2$ from $(-1,1)$ to $(2,4)$ through the vector field $F(x,y) = xe^y\hat{i} + y\hat{j}$.
\end{problem}

\vspace{9cm}

\begin{problem}
    Find the work done by moving along the line segment from $(-1,3,2)$ to $(1,-2,4)$ through the vector field $F(x,y) = (x+2y)\hat{i} +  2z\hat{j} + (x-y)\hat{k}$.
\end{problem}

\newpage

\section{Scratch Work}

\end{document} 